% !TeX root = ../colloquium-2.tex

% Формулировка о разложении по формуле Тейлора с остаточным членом в общей форме.

\phantomsection\label{sec_Definition17}
\begin{definition}[
    Формулировка о разложении по формуле Тейлора с остаточным членом в общей форме.
]
    Пусть в каждой точке отрезка с концами $a$ и $x$ функция $f$ имеет $n$ непрерывных
    производных, а в каждой точке интервала с концами $a$ и $x$ $f$ дифференцируема $n + 1$ раз.
    Пусть также функция $g$ непрерывна на отрезке с концами $a$ и $x$ и дифференцируема в каждой 
    точке интервала с концами $a$ и $x$, причём $g'(\xi) \neq 0$ в любой точке $\xi$ этого 
    интервала. Тогда 
    $$ R_n (x; \; f, \; a) = \frac{f^{(n + 1)} (c)}{n! g'(c)} (g(x) - g(a))(x - c)^n, $$
    где $c$ - какая-то точка из интервала с концами $a$ и $x$.
\end{definition}